% v15. Additional independent module after one_high_and_crossed_pairing_v15.tex.
% The finite calculation concerns only cores 2,3,4 and has no denominator cutoff.

\subsection{二葉三角形の負の枝の排除と残る小 reverse 条件}
\label{tl15:negative-section}

\begin{lemma}[二高分母標準形の $x$ は一でない]
\label{tl15:x-at-least-two}
正の二高分母標準形では $x=(a-1)/r\ge2$ である．
\end{lemma}

\begin{proof}
正の枝 $S>r$ では $x=sS/r>1$ である．
負の枝 $S=T=w<r$ で $x=1$ と仮定する．
$s w=r$ なので $s\ge2$ である．$y=zw/(r+1)$ の整数性と
$\gcd(w,r+1)=1$ から $z=k(r+1)$，$y=kw$，$k\ge1$ と書ける．
\[
 a=r+1,\quad b=k(r+1)w+1,\quad
 A=r+2-s,\quad B=1+k\{r(w+1)+1\}
\]
を $AB=ab-1$ に代入し，$r=sw$ を用いると
\[
 k\{w(-s^2+3s-1)+2-s\}=s-2
\]
となる．$s=2$ なら左辺は $kw>0$，右辺は零である．
$s\ge3$ なら左辺は負，右辺は正である．いずれも矛盾する．
\end{proof}

\begin{lemma}[交差 pairing の大きい reverse entries の障害]
\label{tl15:large-reverse}
定理~\ref{tl15:crossed-pairing} の全仮定の下で，残る交差標識
$(A,B,C,D)=(0,1,4,3)$ を取る．
$r\ge5,S\ge5,z>1$ が同時に成立することはない．
\end{lemma}

\begin{proof}
core $2$ の低 rows は
\[
 (p^2_{30},p^2_{31},p^2_{34})=(-S,T,S-T-1),\qquad
 (p^2_{40},p^2_{41},p^2_{43})=(r,-r-1,0).
\]
各低辺の node 数は非負で総和五なので，各 puncture で
$-4\le p+p'\le1$ である．従って低辺 $23,24$ は
\[
 p^3_{20}\ge S-4>0,\qquad p^4_{21}\ge r-3>0
\]
を要求する．一高分母標準形~\eqref{tl15:leaf-matrix} により，
両葉の ray $2$ はその標準形の $D$ でなければならない．

補題~\ref{tl15:x-at-least-two} と $r>1$ より
$1<x<a-1$ であり，$z>1,T\ge S>1$ より $1<z<b-1$ である．
共有高辺 $02,12$ の共通低 punctures に
補題~\ref{ls11:low-puncture} を適用すると，反対符号の通常の場合しか残らず，
\[
 p^0_{42}=-r,\qquad p^1_{32}=-T
\]
となる．低辺 $04,13$ の node 上界五から，さらに
\[
 p^4_{02}\ge r-4>0,\qquad p^3_{12}\ge T-4>0
\]
を得る．
各葉では ray $2=D$ であり，$B\to D=0$ なので，
これらの正の entries の出発 ray は $C$ である．
従って，両葉で他の葉へ向かう ray は $B$ であり，$c>0,d=0$ となる．
その係数 $D_B=m+t+1+c$ は正なので，低辺 $34$ の両端の係数は相殺できない．
\end{proof}

\begin{lemma}[負の枝で中間 core を固定したときの葉の有限復元]
\label{tl15:finite-leaves}
交差標識 $(A,B,C,D)=(0,1,4,3)$ を持つ負の枝 $S=T=w<r$ の
core $2$ の整数行列を固定する．低辺 $23,24$ において
marked node 条件と $D$ の相殺を満たす一高分母の葉の行列は，
分母の上限を置かず，有限個の約数計算から全て復元できる．
\end{lemma}

\begin{proof}
葉 $3$ で ray $2=B$ とすると $p^3_{21}=0$ であるが，
$p^2_{31}=w>1$ と合わせると低辺 $23$ の node 非負性に反する．
葉 $4$ で ray $2=B$ とすると $p^4_{20}=0$ であるが，
$p^2_{40}=r>1$ と合わせても同じ矛盾を得る．
従って両葉で ray $2$ は $C$ または $D$ である．

葉を $i\in\{3,4\}$，その高 ray を $H_i$ とし，固定された整数を
\[
 R_i=p^2_{iH_i},\qquad E_i=D^2_i
\]
と書く．ray $2=C$ の場合，$p^i_{2H_i}=-u$，$u=th$ は
\[
 u\in[R_i-1,R_i+4]\cap\ZZ_{>0}
\]
に入る．各 $u$ の全正約数 $h$ を取って $t=u/h$ とすれば，
$h,t$ の全 choices を得る．係数相殺は $d=h+E_i$ を固定する．
$d\ne0$ なら $c=0$ であり，$d=0$ なら残る $c=p^i_{2D}$ は
低辺 $2i$ の puncture $D$ の六整数の node 区間に入る．
残る二つの低 ray の $B,D$ の割当は二通りしかない．

ray $2=D$ の場合，$p^i_{2H_i}=v$，$v=t(h+1)$ は
\[
 v\in[-4-R_i,1-R_i]\cap\ZZ_{\ge2}
\]
に入る．各 $v$ の二以上の全正約数 $h+1$ を取って $t=v/(h+1)$ とする．
係数相殺は $c=E_i-h-1$ を固定する．
$c\ne0$ なら $d=0$ であり，$c=0$ なら $d=p^i_{2C}$ が
同じ六整数の node 区間で固定される．$B,C$ の二つの割当を試せばよい．

最後に元の三つの node 条件と係数相殺を再検査すれば，ちょうど全候補を得る．
どの段階にも葉の分母 $th(h+1)+1$ に対する上限は仮定していない．
\end{proof}

\begin{proposition}[小さい負の枝の低三角形は接続しない]
\label{tl15:small-negative}
交差標識で $S=T\le4$ の場合には，cores $2,3,4$ の低辺
$23,24,34$ における node 条件と係数相殺を同時に満たす模型はない．
\end{proposition}

\begin{proof}
補題~\ref{cp14:small-negative} の七模型がこの場合を全て尽くす．
補題~\ref{tl15:finite-leaves} の全約数復元を実行すると，
低辺 $23,24$ にそれぞれ接続する候補数と，検査する葉の全組数は次のとおりである．
\[
\begin{array}{rrrr|rr|r}
 r&s&z&w&\#\text{leaf }3&\#\text{leaf }4&\#\text{pairs}\\\hline
 4&4&5&2&4&8&32\\
 5&10&2&3&4&8&32\\
 6&6&14&3&4&19&76\\
 7&7&2&4&14&10&140\\
 7&21&4&4&5&10&50\\
 8&6&9&4&5&12&60\\
 8&8&27&4&5&19&95
\end{array}
\]
合計四百八十五組は全て，低辺 $34$ の三つの node 区間
$[-4,1]$ または $D^3_4+D^4_3=0$ の少なくとも一方に反する．
標準 library の整数・有理数演算だけを用いる
\nolinkurl{verify_two_leaves_v15.py} は，七模型を生成する完全な約数 trials，
全葉候補の行列，全組の棄却理由を
\nolinkurl{two_leaves_exact_certificate_v15.json} に出力する．
検算器は復元した各行列の元の row sums，pair products，係数，
および正の cycle inequalities を全て再検査する．

これは三つの cores の低三角形だけで生じる矛盾である．
三価 cores の未知の分母 $q_{01}$ を固定・列挙する必要はなく，
その値によってこの矛盾を回避することもできない．
\end{proof}

\begin{theorem}[二葉三角形に残る正の枝と小 reverse 条件]
\label{tl15:remaining-positive}
定理~\ref{tl15:crossed-pairing} の全仮定の下で，二葉三角形 support
$\{01,02,12,03,14\}$ が残るなら，二価 core $2$ は
\[
 (A,B,C,D)=(0,1,4,3),\qquad S>r,\qquad rT>(r+1)S
\]
を満たす正の枝であり，さらに
\begin{equation}\label{tl15:small-reverse-or-unit}
 r\le4\quad\text{または}\quad z=1
\end{equation}
でなければならない．
\end{theorem}

\begin{proof}
非交差標識は定理~\ref{tl15:crossed-pairing} によって排除される．
交差標識で負の枝 $S=T=w$ を考える．$w\le4$ は
命題~\ref{tl15:small-negative} に反する．$w\ge5$ なら
$r>w\ge5$ であり，$y=zw/(r+1)$ は正整数で $y<z$ なので $z>1$ である．
従って補題~\ref{tl15:large-reverse} に反する．
よって正の枝のみが残り，$rT>(r+1)S$ は既知の二分補題に従う．
その枝で $r\ge5,z>1$ なら $S>r\ge5$ なので，
再び補題~\ref{tl15:large-reverse} に反する．
\end{proof}

式~\eqref{tl15:small-reverse-or-unit} は今後の無限族の解析を絞る条件である．
$r\le4$ の範囲でも $S,T,s,z$ に上限を置いたわけではなく，
$z=1$ の枝も残るため，この段階で二葉三角形 support 全体の不存在は結論しない．
